% ===========================================================================
\chapter{Why?}\label{ch:explanations}

The engine knows that Carol and Dan are cousins. A person might reasonably
ask how it knows. The answer is a \emph{proof}: \code{cousin(carol, dan)}
holds by the rule for cousins, because Bob is Carol's parent, Beth is Dan's
parent, and Bob and Beth are siblings; and they are siblings by the rule for
siblings, because Alice is the parent of both. Every step is a rule applied
to facts, down to facts that were simply given.

To give that answer the engine has to remember, for each fact a rule
derives, which rule derived it and with what bindings. Every derived fact
already passes through \code{add_derived} with exactly those two things, so,
once again, a subclass is all it takes.

\section{Remembering derivations}

\begin{newfile}{explain.nex}
-- explain.nex: remembering how each fact was derived, and saying why a
-- fact holds.

intern Database

-- How a fact was derived: by a rule, with bindings that made its body
-- true.
class Derivation
create
  make(r: Rule, b: Bindings) do
    rule := r
    bindings := b
  end
feature
  rule: Rule
  bindings: Bindings
end
\end{newfile}

An explained database is an indexed one that also keeps a map from each
derived fact to its first derivation. The first is enough: one proof is all
we need, and the first one found is never longer than later ones, since
everything it used was known before anything later.

\begin{continued}{explain.nex}
-- A database that remembers, for every fact a rule derived, the first
-- derivation it found.
class Explained_Database inherit Indexed_Database
create
  make do
    Indexed_Database.make
    how := {}
  end
feature
  how: Map[String, Derivation]

  add_derived(a: Atom, rule: Rule, b: Bindings): Boolean do
    result := add(a)
    if result then
      how.put(a.text, create Derivation.make(rule, b))
    end
  ensure
    remembered: how.contains_key(a.text) or not result
  end

  is_derived(a: Atom): Boolean do
    result := how.contains_key(a.text)
  end

  -- The proof that `a` holds, one step per line: each derived fact,
  -- with the rule that gave it, above the facts that rule used.
  why(a: Atom): Array[String]
  require
    known: contains(a)
  do
    result := []
    explain(a, "", result)
  ensure
    every_step_checks: steps_hold(a)
  end
\end{continued}

\section{Printing a proof}

A proof prints from the top down: the fact and the rule that derived it,
then, indented below, each condition of the rule's body with the bindings
applied. An atom in the body is a fact itself, so it is explained the same
way, recursively, until it reaches a fact that was given. A test is shown
as it was checked, and a negated atom as the fact that does not hold.

\begin{continued}{explain.nex}
  explain(a: Atom, indent: String, lines: Array[String]) do
    if is_derived(a) then
      let d := how.get(a.text)
      lines.add(indent + a.text + "    by " + d.rule.text)
      across d.rule.body as l do
        match l of
          Positive(atom) then explain(substitute(atom, d.bindings), indent + "  ", lines)
          Compare(op, left, right) then
            lines.add(indent + "  " + literal_text(substitute_literal(l, d.bindings)))
          Negated(atom) then
            lines.add(indent + "  " + literal_text(substitute_literal(l, d.bindings)) + "    (no such fact)")
        end
      end
    else
      lines.add(indent + a.text + "    (given)")
    end
  end
\end{continued}

Does the recursion always end? Each step goes from a fact to facts that were
in the database \emph{before} it was derived, since the join that derived it
read them there. Going back through ``derived before'' cannot go round in a
circle, so it ends.

\section{Checking a proof}

A proof is only worth printing if it is right, and \code{why}'s
postcondition checks that it is: every step's rule, with its bindings,
really does give that fact, every atom it used is in the database, every
test holds, and every negated atom is absent. If a bug ever recorded the
wrong bindings for a fact, the explanation would be caught before anyone
read it.

\begin{continued}{explain.nex}
  -- Is the recorded derivation of `a`, and of every fact it used, a true
  -- instance of its rule in this database?
  steps_hold(a: Atom): Boolean do
    result := true
    if is_derived(a) then
      let d := how.get(a.text)
      result := substitute(d.rule.head, d.bindings).text = a.text
      across d.rule.body as l do
        match l of
          Positive(atom) then
            result := result and contains(substitute(atom, d.bindings))
                      and steps_hold(substitute(atom, d.bindings))
          Compare(op, left, right) then
            result := result and holds(op, substitute_term(left, d.bindings),
                                       substitute_term(right, d.bindings))
          Negated(atom) then
            result := result and not contains(substitute(atom, d.bindings))
        end
      end
    end
  end
end
\end{continued}

\begin{continued}{explain.nex}
-- `l` with its bound variables replaced by their values.
function substitute_literal(l: Literal, b: Bindings): Literal do
  match l of
    Positive(atom) then result := create Positive.make(substitute(atom, b))
    Compare(op, left, right) then
      result := create Compare.make(op, substitute_term(left, b), substitute_term(right, b))
    Negated(atom) then result := create Negated.make(substitute(atom, b))
  end
end
\end{continued}

\section{Trying it out}

\begin{tryit}{try_why.nex}
intern Parser
intern Analysis
intern Explain
intern io/Path

match parse_program(create Path.make(program.working_directory + "/family.dl").read_text) of
  Ok as ok then
    do
      let db := create Explained_Database.make
      evaluate_program(ok.value, db, create Semi_Naive_Evaluator.make)
      match parse_program("?- ancestor(alice, emma).") of
        Ok as q then
          across db.why(q.value.queries.get(0)) as line do
            console.print_line(line)
          end
        Err as e then console.print_line(e.error)
      end
    end
  Err as e then console.print_line(e.error)
end
\end{tryit}
\tryoutput{try_why}
